\subsection{Insights and Design Implications}
\label{sec:design-insights}

Three premises connect workflow structure to selective KV protection:
estimating an agent's remaining time, identifying workloads with sparse
historical reads, and predicting the reads of an uncommitted retry tail.
Figure~\ref{fig:design-insights} separates their evidence. The first
premise remains a validation gap in this draft; the other two motivate
the placement policy within the scope stated below.

\begin{figure*}[t]
  \centering
  \captionsetup[subfigure]{font=normalsize,labelfont=normalfont,
    textfont=normalfont,justification=centering,singlelinecheck=false}
  \begin{subfigure}[t]{0.32\textwidth}
    \vspace{0pt}
    \includegraphics[width=\linewidth]{figures/design-insights/agent-time-audit.pdf}
    \caption{Agent-time prediction}
    \label{fig:insight-agent-time}
  \end{subfigure}\hfill
  \begin{subfigure}[t]{0.32\textwidth}
    \vspace{0pt}
    \includegraphics[width=\linewidth]{figures/design-insights/sparse-attention-survey.pdf}
    \caption{Sparse attention adoption}
    \label{fig:insight-sparse-adoption}
  \end{subfigure}\hfill
  \begin{subfigure}[t]{0.32\textwidth}
    \vspace{0pt}
    \includegraphics[width=\linewidth]{figures/design-insights/retry-kv-prediction.pdf}
    \caption{Retry KV prediction}
    \label{fig:insight-retry-kv}
  \end{subfigure}
  \caption{Evidence for the three design premises.
    (a) Offline five-fold prediction audit of 1,688 CC agents from 174
    workflows; workflows do not cross folds. Each agent contributes its
    mean absolute remaining-time error divided by total duration.
    The symmetric-log axis retains all errors; all recorded roles are
    \texttt{Subagent}. This does not establish accurate role-conditioned
    timing prediction.
    (b) Four representative 2025--2026 releases, their sparse mechanisms,
    and reported context limits; DS denotes DeepSeek. MiniCPM4 uses its
    native 32K limit. This is a source survey, not an adoption-rate estimate.
    (c) Native PG19 traces, 32K history, three documents per model.
    Past-frequency ranking observes queries 32,257--32,768; demand is
    the distinct historical KV selected by queries 32,769--32,832.
    Both policies reserve mandatory blocks within the same 40\% budget.
    Bars are document means; whiskers span the three documents. This is
    trace-derived hypothetical retry demand, not a fault-recovery run.}
  \label{fig:design-insights}
  \Description{Three side-by-side panels show a negative agent timing
    prediction audit, a primary-source survey of sparse model releases,
    and held-out retry KV coverage. Past-frequency ranking improves
    coverage substantially for NOSA and only modestly for MiniCPM4.}
\end{figure*}

\paragraph{Insight 1: workflow-aware protection requires validated
remaining-time estimates.}
Repeated agent roles offer a basis for profiling: prior work observes
role-dependent output-length distributions~\cite{yu2026pythia}.
Output length, however, is not a multi-round agent's wall-clock duration,
which also includes tool execution and service delays. We therefore
audit remaining-time prediction on completed CC agents~\cite{semianalysis2026cctraces},
using only features available at each request arrival and withholding
whole workflows. A diagnostic progress GBDT has a median normalized
error of 51.8\%, versus 49.3\% for a training-fold median prior
(Figure~\ref{fig:design-insights}\subref{fig:insight-agent-time}).
The trace labels every agent as \texttt{Subagent}, preventing validation
of the role-specific P90 estimator used by the Tracker.

\emph{Design implication.} AegisKV separates time estimation from DAG
propagation and conservatively handles unknown successors
(Section~\ref{sec:workflow-protection}). We still need repeated timings
for each role, interval coverage, and protection tier error rates.
We do not treat Figure~\ref{fig:insight-agent-time} as
positive evidence of accurate agent-time prediction or replace the
current P90 policy with the diagnostic GBDT.

\paragraph{Insight 2: sparse attention is relevant to long-context serving.}
MiniCPM4 uses InfLLMv2~\cite{minicpm4model}; DeepSeek-V3.2 and GLM-5
use DeepSeek Sparse Attention and explicitly target agentic
capabilities~\cite{deepseek2025v32model,zai2026glm5model}.
DeepSeek-V4 combines compressed sparse and heavily compressed attention
at a reported million-token context~\cite{deepseek2026v4model}.
Figure~\ref{fig:design-insights}\subref{fig:insight-sparse-adoption}
documents these examples; it does not establish that most deployments
already use sparse attention.

\emph{Design implication.} AegisKV exposes actual layer/block selections
to the Planner instead of assigning equal value to every historical
KV byte. This interface must preserve selector state and all required
attention branches. Different sparse or compressed models need their
own recovery objects and calibration; the survey does not demonstrate
AegisKV compatibility with every listed model.

\paragraph{Insight 3: past selections predict part of a short retry
tail's distinct historical KV demand.}
After failure, rebuilding an uncommitted suffix replays known token IDs.
The relevant forecast is which \emph{historical} blocks these queries
will need, before their own selector outputs are available. On native
NOSA-1B, NOSA-3B, and MiniCPM4-8B traces, we use a strictly earlier
512-query window to rank historical blocks and hold out the next
64 queries. At the same 40\% historical-byte budget, past frequency
covers 79.2\%, 77.8\%, and 53.6\% of distinct needed KV, versus
59.7\%, 59.7\%, and 50.9\% for recency
(Figure~\ref{fig:design-insights}\subref{fig:insight-retry-kv}).
The improvement is model-dependent and does not recover every needed
block. These are PG19 demand traces, not agentic failure-retry executions.

\emph{Design implication.} The Planner uses past selections for advance
placement; exact demand fetches cover misses (Section~\ref{sec:kv-placement}).
Repeated reads of one block count once when sizing recovery bytes.
With recoverable selector state and the same causal prefix and token
IDs, the held-out demand is a proxy for retry demand; physical replay
equivalence still requires validation.

Together, these premises motivate workflow-level protection tiers and
request-internal block selection. The remaining design problem is to
turn them into a capacity-bounded plan that removes exposed transfer
wait, then execute independent retries without imposing a serial
recovery order (Sections~\ref{sec:kv-placement}--\ref{sec:kv-runtime}).
